The tools available to people working in artificial intelligence change
quickly. An advanced course cannot be limited to learning whichever model or
framework is currently in use. In Advances in Artificial Intelligence,
part of the DAISY programme at Hochschule Düsseldorf, I want students to
understand how a method works, what assumptions it makes, and what evidence
would justify choosing it for a particular problem. Deep learning is an
important part of that discussion, but it is not a default answer. Classical
methods, domain knowledge, and newer AI techniques can be combined when the
problem calls for them. The ability to judge those choices will remain useful
longer than familiarity with any one tool.

The course project gives students room to put that judgment into practice.
Each team develops a topic with me and then carries much of the work forward
independently over the semester. Regular conversations provide opportunities
to question the scope, examine technical decisions, and respond to results
that do not match initial expectations. Students must decide what can be
implemented in the available time, learn methods beyond those covered in
class, and explain why their particular design is appropriate. The amount of
code written is not, by itself, a measure of achievement. With AI coding
tools now part of many development workflows, it matters even more that
students can explain their own implementation in detail and understand the
consequences of changing it.

The projects culminate in the KI-Con, an annual public event that is
closer to a small conference or exhibition than to a sequence of final
presentations. Each team prepares an A1 poster, a working live demonstration,
and the opportunity for visitors to ask questions and interact with the
project. The poster should make the problem, approach, and findings accessible
at a glance. The demonstration should make the implemented system tangible.
Conversation then gives students a chance to explain their reasoning, defend
their choices, and acknowledge what remains uncertain. These exchanges are
part of the educational value of the event, not merely a way to display a
finished product.

This volume adds a third perspective. Each team has written a short
scientific paper, limited to four pages, to document the question it pursued,
the relevant prior work, its method, the evidence it gathered, and the limits
of its conclusions. A live demonstration shows what was built; a poster gives
a compact visual account; a paper allows the argument and its supporting
details to remain available after the day of the KI-Con. Bringing the papers
together makes it possible to read the projects as a record of the questions
and decisions that shaped this year's work, rather than only as a collection
of final outputs.

Not every project will yield a complete system or an unambiguous success.
Unexpected results, approaches that prove unsuitable, and decisions that a
team would now make differently are normal parts of research and development.
Recognizing and explaining these limits is itself evidence of serious work.
I thank the students for the care they have put into their projects and for
their willingness to discuss both what succeeded and what did not. I hope
these proceedings serve as a useful snapshot of the questions that students
of data science and artificial intelligence chose to investigate at this
point in time.

